\section{Graph reductions and homology}\label{sec:prelim}

The graph operations below are determined by induced neighbourhoods.
For a vertex $v$ of $G$, we write $N_G(v)$ for the subgraph induced by
its neighbours. For distinct vertices $u,v$, the common neighbourhood
$N_G(u,v)$ is induced by the vertices adjacent to both. We write $G-v$
for a vertex deletion, $G-uv$ for an edge deletion, and $G+uv$ for the
addition of a missing edge. The clique complex $\cl(G)$ consists of the
cliques of $G$; a complex of this form is called \defi{flag}.
The clique complexes of $N_G(v)$ and $N_G(u,v)$ are the links of $v$
and, when $uv$ is an edge, of $uv$ in $\cl(G)$.

\subsection{Contractible and strong contractible graphs}
\label{sec:contractible}
The two classes differ in the permitted operations and in the
condition imposed on neighbourhoods.

\begin{definition}[{\cite[Definition~1]{DEFH2023}}]\label{def:I}
The class $\I$ of \defi{I-contractible graphs} is the least class
containing $K_1$ and closed under the following operations:
\begin{enumerate}[label=(\roman*),leftmargin=2.5em]
\item delete a vertex $v$ when $N_G(v)\in\I$;
\item add a new vertex whose neighbourhood is an induced subgraph
of $G$ belonging to $\I$;
\item delete an edge $uv$ when $N_G(u,v)\in\I$;
\item add a missing edge $uv$ when $N_G(u,v)\in\I$.
\end{enumerate}
The empty graph does not belong to $\I$.
\end{definition}

The operations in Definition~\ref{def:I} preserve the simple homotopy
type of the clique complex \cite[Theorem~2.8]{DochtermannMatsushita2025};
the graph formulation goes back to Ivashchenko and to Chen, Yau and Yeh
\cite{Ivashchenko1994a,ChenYauYeh2001}. In particular, every member
of $\I$ has contractible clique complex. The converse holds as well
\cite[Corollary~1.5]{DochtermannMatsushita2025}.

\begin{definition}[{\cite[Definition~3 and Lemma~1]{EFH2022};
\cite[Definitions~4--5]{DEFH2023}}]\label{def:IS}
The class $\IS$ of \defi{strong I-contractible graphs} contains $K_1$
and is defined recursively, for graphs of order at least two, by
\[
 G\in\IS\quad\Longleftrightarrow\quad
 \begin{cases}
 G-v,\ N_G(v)\in\IS&\text{for some }v\in V(G),\quad\text{or}\\
 G-uv,\ N_G(u,v)\in\IS&\text{for some }uv\in E(G).
 \end{cases}
\]
The class $\ISV$ of \defi{strong vertex I-contractible graphs} has the
same base case and uses only the first alternative, with both smaller
graphs required to belong to $\ISV$. Neither class contains the empty
graph.
\end{definition}

Each recursive argument has smaller $|V|+|E|$, so these definitions are
well founded. Reversing a successful reduction gives the equivalent
definitions by gluings in the cited papers. Induction on $|V|+|E|$
gives
\begin{equation}\label{eq:strong-inclusions}
 \ISV\subseteq\IS\subseteq\I.
\end{equation}
For the second inclusion, the remaining graph and the tested
neighbourhood lie in $\I$ by induction, and the inverse operation
reconstructs the original graph by an allowed gluing.

An operation is \defi{$\IS$-admissible} if its tested neighbourhood
belongs to $\IS$. Unless a different class is named, admissibility
will mean $\IS$-admissibility. A \defi{strong reduction} is a sequence
of admissible deletions ending at $K_1$. The empty-neighbourhood
convention means that neither an isolated vertex nor an edge with empty
common neighbourhood can be deleted. Every nonempty tree belongs to
$\ISV$: deleting a leaf tests the neighbourhood $K_1$ and reduces the
order. Conversely, a triangle-free member of $\I$ is a tree, since
its clique complex is connected, contractible and of dimension at most one.

Membership in $\IS$ requires the existence of a successful reduction.
A different admissible first deletion can destroy that possibility,
as the examples for Conjecture~15 will show. Thus the recursive test
must consider all possible first deletions.

\subsection{Dismantling levels}
Vertex reductions also give a hierarchy in which the condition on the
neighbourhood changes with the level.

\begin{definition}[{\cite[Definition~1]{FieuxJouve2020}}]\label{def:kdis}
A \defi{cone} is a nonempty graph with a vertex adjacent to every other
vertex. Cones form the class $(-1)$-Dis. For $k\geq0$, a vertex $v$
is \defi{$k$-dismantlable} when $N_G(v)$ belongs to $(k-1)$-Dis.
A graph is \defi{$k$-dismantlable}, or belongs to $k$-Dis, if it can
be reduced to $K_1$ by successive deletions of $k$-dismantlable
vertices. The singleton belongs to every level, and the empty graph
to none.
\end{definition}

At level zero, the condition says that some neighbour $u$ of $v$ is
adjacent to all the other neighbours of $v$. Equivalently,
$N_G[v]\subseteq N_G[u]$, where brackets denote closed neighbourhoods.
The hierarchy satisfies
\begin{equation}\label{eq:dismantling-inclusions}
 (-1)\text{-Dis}\subseteq0\text{-Dis}\subseteq1\text{-Dis}
 \subseteq\cdots\subseteq\ISV\subseteq\IS\subseteq\I;
\end{equation}
see \cite[Proposition~2]{FieuxJouve2020} and
\cite[Lemma~19 and Theorem~23]{DEFH2023}.
The union of the finite dismantling levels is $\ISV$
\cite[Theorem~23]{DEFH2023}.

\subsection{Coefficient choices}
We use homology to distinguish graphs whose clique complexes are
contractible from graphs with the same homology over $\mathbb F_2$.
Write $\widetilde H_i(G;R)=\widetilde H_i(\cl(G);R)$.
For $R=\ZZ$ or $\mathbb F_2$, let $\mathcal A_R$ be the class of
nonempty graphs for which $\widetilde H_i(G;R)=0$ for every $i\geq0$.
The Euler characteristic is
\[
 \chi(G)=\sum_{i\geq0}(-1)^if_i(G),
\]
where $f_i(G)$ counts the cliques with $i+1$ vertices.

For a nonempty finite complex $K$, the universal coefficient theorem
\cite[Theorem~3A.3]{Hatcher2002} gives, for $i\geq0$,
\begin{equation}\label{eq:uct}
 0\longrightarrow\widetilde H_i(K;\ZZ)\otimes\mathbb F_2
 \longrightarrow\widetilde H_i(K;\mathbb F_2)
 \longrightarrow\operatorname{Tor}^{\ZZ}_1
       (\widetilde H_{i-1}(K;\ZZ),\mathbb F_2)
 \longrightarrow0.
\end{equation}
We use the augmented convention $\widetilde H_{-1}(K;\ZZ)=0$.
Consequently,
\begin{equation}\label{eq:homology-inclusions}
 \IS\subseteq\I\subseteq\mathcal A_{\ZZ}
 \subseteq\mathcal A_{\mathbb F_2}.
\end{equation}
For a connected finite complex, $\mathbb F_2$-acyclicity allows finite
integral homology of odd order. In particular, both
$(\ZZ/3)\otimes\mathbb F_2$ and
$\operatorname{Tor}^{\ZZ}_1(\ZZ/3,\mathbb F_2)$ vanish.
This accounts for the different hypotheses in the two assertions of
Conjecture~14. Throughout the article, unqualified acyclicity means
membership in $\mathcal A_{\mathbb F_2}$.

\subsection{A disc with boundary identifications}\label{sec:discs}
The first construction is most easily understood by obtaining its
clique complex from a triangulated disc. The following criterion
determines the homotopy type from the boundary identifications.

\begin{proposition}\label{prop:disc}
Let $D$ be a triangulated closed disc and $q:D\to X$ a simplicial
map. Suppose that $q$ is injective on the interior of $D$ and identifies
$X$ with the quotient of $D$ obtained solely by identifying boundary
points having the same image. If $S=q(\partial D)$ is a connected
graph with exactly one cycle, then
\[
 X\simeq S^1\cup_d D^2,
\]
where $d$ is the degree of the boundary map after contracting a maximal
tree in $S$. For $|d|=1$ the complex is contractible. For $|d|=3$
it is homotopy equivalent to $M(\ZZ/3,1)$ and has
$\widetilde H_1(X;\ZZ)\cong\ZZ/3$ and zero reduced homology in all
other degrees.
\end{proposition}
\begin{proof}
The hypotheses identify $X$ with the adjunction space $S\cup_fD^2$.
Collapsing a maximal tree $T$ in $S$ is a homotopy equivalence, since
$T$ is a contractible subcomplex \cite[Propositions~0.16--0.17]{Hatcher2002}.
The quotient is $(S/T)\cup_{pf}D^2$, with $S/T\cong S^1$.
Maps $S^1\to S^1$ are classified up to homotopy by their degree
\cite[Theorem~1.7]{Hatcher2002}, and homotopic attaching maps have
homotopy equivalent adjunction spaces
\cite[Proposition~0.18]{Hatcher2002}. This proves the asserted
homotopy type. For $|d|=1$ the map is homotopic to a homeomorphism,
whose mapping cone is contractible. For $|d|=3$ this is the defining
adjunction space of $M(\ZZ/3,1)$, and its cellular chain complex in
positive degrees is $0\to\ZZ\xrightarrow{d}\ZZ\to0$.
\end{proof}
